\chapter{\rm{VASP}的基本使用:~输入文件与控制参数}
用户使用\textrm{VASP}，通常需要准备四个输入文件，%\textrm{VASP}的计算输入文件主要
包括\textrm{POSCAR}、\textrm{INCAR}、\textrm{KPOINTS}和\textrm{POTCAR}。\footnote{\textrm{VASP}中的文件名，很多是以\textrm{``CAR''}结尾命名的，这源自德语\textrm{``die Karte''}。这个词有``记录、档案''的意思。}
这四个文件分别是计算对象的结构文件，计算控制参数文件，$\vec k$空间布点和积分方案文件和用于计算的元素的\textrm{PAW}方法所涉及的原子数据。
这四个文件中，\textrm{POTCAR}的数据是\textrm{VASP}软件包自带的。其余三个文件，则需要用户按一定的规则输入或构建。这四个文件必须放置在同一工作目录下，且文件名必须全部大写、无扩展名，文件内容必须是纯文本格式。%这四个文件各自承担着不同的职责——INCAR是“控制中心”，POSCAR是“结构蓝图”，描述原子体系的几何构型；POTCAR是“能量密码本”，包含各元素的赝势信息；KPOINTS是“采样方案”，定义布里渊区中k点的取样方式。
四个文件的内容共同决定了计算的精度、效率和最终结果的可靠性。


\section{\rm{POSCAR}文件与晶胞参数}
\textrm{POSCAR}文件保存的是计算对象的结构文件，包括晶胞几何结构和元胞中原子(或离子)的位置，在分子动力学计算中，还有每个原子的初始速度信息。
\subsection{\rm{POSCAR}文件的内容}
{\fontsize{9.2pt}{7.2pt}\selectfont{
\verbatiminput{Figures/POSCAR.txt} %为保险:~选用文件名绝对路径
}}
\textrm{POSCAR}文件的格式基本是固定的。第一行字符串，描述的是计算对象的名称；第二行对应的是晶胞参数的通用标度系数(如果这个行的值是负值，则表示晶胞体积)；后续三行(第三到第五行)以矩阵形式给定晶胞参数$a$、$b$、$c$的矢量方向(等价于确定晶胞参数的$\alpha$、$\beta$和$\gamma$)。第六行给定计算体系各原子的原子数，此处的原子顺序将于\textrm{POTCAR}中原子顺利一致。如果是分子动力学或结构计算时，第七行必须是特征字符串\textrm{Selective dynamic}(至少也要是保持首字符\textrm{S}或\textrm{s})。\textrm{VASP}程序读到该关键字符串时，可以要求每个原子沿晶胞三个轴向是否需要移动。接下来一行(如果设定\textrm{Selective dynamic}时是第八行，否则是第七行)，将给出晶胞中各原子坐标的格式，许可的特征字符串是\textrm{Caresian}(对应笛卡尔坐标)或\textrm{Dirct}(分数坐标)。类似地，\textrm{VASP}只依据字符串首字母确定坐标特征，\textrm{C}，\textrm{c}或\textrm{K}或\textrm{k}是笛卡尔坐标，\textrm{D}，\textrm{d}表示分数坐标。接下来各行是对应原子在晶胞中的坐标位置。在分数坐标系下，每个原子的位置可以表示为
\begin{equation}
	\vec R=x_1\vec a_1+x_2\vec a_2+x_3\vec a_3
	\label{eq:diract-mode}
\end{equation}
其中$a_{1\cdots 3}$是晶胞的三个基矢，$x_{1\cdots 3}$即各行依次给出的数值。而笛卡尔坐标系下，每个原子的位置则须乘上第二行的通用标度系数
\begin{equation}
	\vec R=s
	\begin{pmatrix}
		\begin{aligned}
			x_1\\x_2\\x_3
		\end{aligned}
	\end{pmatrix}
	\label{eq:Cartensian-mode}
\end{equation}
注意，各原子的排列顺序和数目必须与第六行的原子顺序与对应的数目保持一致。而一旦第七行是\textrm{Selective dynamic}时，每一行除了原子坐标值外，还可以通过额外的三个字符\textrm{T}或\textrm{F}标记原子在三个方向的运动或固定。
%\subsection{晶体与晶胞参数}
%\centering
%\vspace{-0.15in}
%------------------------------------直-接-插-入-文-件--------------------------------------------------------------------------------------
%\textcolor{red}{\textbf{直接插入文件}}:
%\fontsize{4.2pt}{3.2pt}\selectfont{
%\verbatiminput{Figures/POSCAR} %为保险:~选用文件名绝对路径
%}

\section{\rm{INCAR}文件}
\textrm{INCAR}是\textrm{VASP}计算过程中最核心的控制文件，它决定了计算的内容和如何算；\textrm{INCAR}文件中通过一系列``标签\textrm{(tag)} = 值\textrm{(value)}''的键值模式来实现计算参数的控制。
\textrm{INCAR}文件允许包含的参数有数十个，它们决定了计算的任务类型(如结构优化、静态计算、分子动力学等)、电子结构求解算法、收敛标准、输出控制等方方面面。
\textrm{INCAR}的参数虽然繁多，但幸运的是，\textrm{VASP}为\textrm{INCAR}的所有控制参数都设定了可缺省的默认值。因此对于初学者来说，\textrm{INCAR}即使是一个空文件(但这个文件必须存在！)，同样可以完成一个\textrm{VASP}计算。但是因为\textrm{INCAR}涉及大量参数，这种轻率的做法往往是\textrm{VASP}计算中最常见的错误来源。事实上，用户要想完成一个精准的\textrm{VASP}计算，\textrm{INCAR}中的参数往往不能仅依赖默认的缺省值。
%VASP为每个标签都提供了合理的默认值，因此用户不需要在\textrm{INCAR}中指定所有参数——用户可以不指定任一参数，全部采用默认值。
对\textrm{VASP}用户来说，学习并掌握\textrm{INCAR}中参数的设计与选择，是精通\textrm{VASP}计算的主要任务之一。

%\textrm{POSCAR}、\textrm{KPOINTS}和\textrm{POTCAR}都属于极易处理的输入控制文件，
\subsection{\rm{INCAR}文件的主要内容}
\textrm{INCAR}是一个标签(控制参数)格式的自由\textrm{ASCII}文本文件，其基本语法规则如下
\begin{itemize}
	\item 基本格式:~每行通常为一个``\textrm{tag} = \textrm{value}''的语句\\
		例如:~\textrm{ENCUT = 520}，注意，\textrm{tag}名对大小写敏感，建议严格按照\textrm{VASP}手册规定书写
	\item 多参数同行:~多个参数放在同一行时，彼此用分号\textrm{(;)}间隔\\
		例如:~\textrm{ISMEAR = -5 ; SIGMA = 0.05}
	\item 续行:~可以使用反斜杠\textrm{$\textbackslash$}进行多行续写
	\item 注释:~注释部分必须使用 \# 或! 引领
	\item 空行:~\textrm{INCAR}文件中可以存在空行，通常在读入时会被忽略
\end{itemize}

\textrm{INCAR}的主要参数列于表
\begin{longtable}[c]{|@{\extracolsep{\fill}}c|@{\extracolsep{\fill}}l|}
%\begin{table}[h!]
%    \centering
%    \begin{tabular*}{\textwidth}{|@{\extracolsep{\fill}}c|@{\extracolsep{\fill}}l|}
        \toprule
	\textrm{NGX, NGY, NGZ} &轨道波函数\textrm{FFT}用的网格布点(粗网格)\\
	\textrm{NGXF, NGYF, NGZF} &电荷密度\textrm{FFT}用的网格布点(密网格)\\
	\textrm{NBANDS} &计算的能带数目\\
%	\textrm{NBLK} &\textrm{blocking for some BLAS callls}\\
        \midrule
	\textrm{SYSTEM} &\makecell[tl]{计算对象的名称字符串，仅用于标识计算体系，\\不影响计算结果}\\
%	\textrm{NWRITE} &指定更详细的输出内容到\textrm{OUTCAR}\\
	\textrm{ISTART} &\makecell[tl]{指定初始波函数来源:~0-从头计算; 1-从\textrm{WAVECAR}\\文件读取; 2-计算体积-能量曲线时使用}\\
	\textrm{ICHARG} &\makecell[tl]{指定电荷密度来源:~1-从\textrm{CHARGCAR}文件读取; \\2-由原子电荷密度叠加生成}\\
	\textrm{ISPIN} &是否打开自旋极化:~1-不打开; 2-打开\\
	\textrm{MAGMON} & 指定每个原子上的初始磁矩\\
	\textrm{INIWAV} &\makecell[tl]{针对\textrm{ISTART = 0}时，初猜波函数的生成方案:\\0-简易方案; 1-轨道电子占据数采用随机数}\\
	\midrule
	\textrm{ENCUT} &计算中的能量截断(单位:~\textrm{eV})\\
	\textrm{PREC} &计算精度控制:~\textrm{high/medium/low/normal/accurate}\\
	\textrm{NELM, NELMIN, NELMDL}:~ &指定电子迭代的步数控制参数\\
	\textrm{EDIFF}:~ &指定电子迭代的收敛标准\\
	\textrm{EDIFFG}:~ &指定离子运动的收敛标准\\
	\midrule
	\textrm{NSW} & 指定离子运动的总步数控制参数\\
%	\textrm{NBLOCK/KBLOCK} &\textrm{inner block; outer block}\\
%	\textrm{SYMPREC} &指定对称性精度控制参数\\
	\textrm{LCORR} &受力计算中加入\textrm{Harris}校正\\
	\textrm{POTIM} &指定离子运动模拟中的时间步长(单位:~\textrm{fs})\\
	\textrm{TEBEG, TEEND} &指定原子运动模拟的温度\\
	\textrm{SMASS} &恒温\textrm{Nose}方法的质量参数(单位:~\textrm{am})\\
%	\textrm{NPACO, APACO} &对相关函数\\
	\midrule
	\textrm{POMASS} &用原子单位表示的离子质量\\
	\textrm{ZVAL} &离子电荷\\
	\textrm{RWIGS} &原子的\textrm{Wigner-Sietz}元胞半径\\
	\textrm{NELECT} &体系的总电荷数\\
	\textrm{NUPDOWN} &将体系的自旋磁矩指定为固定值\\
	\textrm{EMIN, EMAX} &指定态密度文件\textrm{DOSCAR}的能量上下限\\
	\textrm{ISMEAR} &\makecell[tl]{指定$\vec k$空间和\textrm{Fermi}面积分方案:~0-\textrm{Gaussian}积分; \\>0-\textrm{MP}方案; -4-四面体积分方案; \\-5-\textrm{Bl\"ochl}改进的四面体积分方案}\\
	\textrm{SIGMA} &\makecell[tl]{指定\textrm{Fermi}面积分采用的展宽系数(单位ie:~\textrm{eV}):\\-4-四面体积分方法; -1-\textrm{Fermi}分布函数; \\0-\textrm{Gaussian}分布函数}\\
	\midrule
	\textrm{ALGO} &\makecell[tl]{指定电子步矩阵对角化算法:\\\textrm{Normal(Davidson)-Fast-Very\_fast(RMM-DIIS)}}\\
	\textrm{IALGO} &\makecell[tl]{指定离子步收敛算法:~8表示共轭梯度\\\textrm{(Conjugate Gradient, CG)}法，\\48表示\textrm{RMM-DIIS}}\\
%	\textrm{RPOT} &指定实空间的非局域投影的网格点数目\\
	\textrm{GGA} &指定交换-相关泛函形式:~可选的如\textrm{PE},\textrm{AM},\textrm{91}等\\
	\textrm{VOSKOWN} &指定\textrm{VWN(Vosko-Wilk-Nusair)}泛函的插值方案\\
	\textrm{DIPOL} &为偶极计算指定元胞中心位置\\
	\textrm{AMIX, BMIX} &电子步迭代时，指定电荷密度混合参数\\
%	\textrm{WEIMIN, EBREAK, DEPER} &指定特定的控制参数\\
%	\textrm{TIME} &指定特定的控制参数\\
	\midrule
	\textrm{LWAVE,LCHARG,LVTOT,LVHAR} &指定输出对应数据到文件\textrm{WAVECAR/CHGCAR/LOCPOT}\\
	\textrm{LELF} &指定输出对应数据到文件\textrm{ELFCAR}\\
	\textrm{LORBIT} &指定输出对应数据到文件\textrm{PROOUT}\\
	\midrule
	\textrm{NPAR} &指定并行计算中，将能带按指定参数分配\\
%	\textrm{LSCALAPACK} &关闭并行计算的\textrm{scaLAPACK}\\
	\textrm{LSCALU} &关闭矩阵对角化的\textrm{LU}分解\\
%	\textrm{LASYNC} &指定并行计算中的重叠通信\\
        \bottomrule
%    \end{tabular*}
    \caption{\textrm{INCAR}常用参数的简要描述}
\label{Tab:INCAR-Tag-0}
%\end{table}
\end{longtable}
\subsection{\rm{VASP}的主要计算控制参数}
%# VASP四大输入文件详解：INCAR、POSCAR、POTCAR与KPOINTS
%
%## 引言
%VASP（Vienna Ab initio Simulation Package，维也纳第一性原理模拟软件包）是一个基于Fortran编写的从头计算软件包，主要用于密度泛函理论（DFT）计算，以处理材料的电子结构问题。它采用平面波基组和赝势方法，适用于模拟固体、表面、分子和纳米结构的电子、结构和动力学性质。VASP的主要功能涵盖电子结构计算、离子弛豫（结构优化）、分子动力学模拟、能带结构分析以及杂化泛函、GW方法和机器学习力场等高级功能。经过数十年的发展，VASP已成为材料科学、凝聚态物理和计算化学领域的标准工具之一。
%要成功运行一次VASP计算，通常需要准备四个核心输入文件：**INCAR**、**POSCAR**、**POTCAR**和**KPOINTS**。
%本文将全面、深入地逐一解析这四大输入文件，从文件格式、关键参数到实际应用技巧，力求为读者提供一份系统而详尽的参考指南。
%## 一、INCAR文件：VASP计算的“控制中心”
%### 1.1 INCAR文件概述
%### 1.2 INCAR文件格式
%### 1.3 INCAR关键参数详解
%
\textrm{INCAR}中的参数众多，按功能类别，对最常用和最重要的参数进行必要的解析和说明。

\subsubsection{系统描述与初始化参数}
\textrm{SYSTEM}:~例如:~\textrm{SYSTEM = Si bulk relaxation}。

\textrm{ISTART}:~控制是否从已存在的\textrm{WAVECAR}文件中读取波函数数据。
\begin{itemize}
	\item 0:~从头开始计算(默认值，一般用于新任务)
	\item 1:~从\textrm{WAVECAR}中读取数据，用于续算，方便加速收敛
	\item 2:~在计算体积-能量(状态方程)相关曲线时使用
	\item 3:~在重启分子动力学计算时使用，需同时存在\textrm{WAVECAR}和\textrm{CHGCAR}文件
\end{itemize}

\textrm{ICHARG}:~控制电荷密度初始化的方式
\begin{itemize}
	\item 0:~由\textrm{WAVECAR}计算电荷密度~(若存在\textrm{WAVECAR})
	\item 1:~从\textrm{CHGCAR}文件读取电荷密度
	\item 2:~由程序通过初猜产生(默认值，一般用于新计算)
	\item 11/12/13:~用于非自洽计算(如\textrm{DOS}和能带计算)
\end{itemize}

\subsubsection{电子结构计算参数}
\textrm{ENCUT}:~平面波截断能量(单位:~\textrm{eV})，定义平面波基组展开的能量截断上限。\textrm{ENCUT}是影响计算精度和效率的重要参数。
\begin{itemize}
	\item \textrm{POTCAR}中\textrm{ENMAX}的值，是确定\textrm{ENCUT}下限的重要参考
	\item \textrm{ENCUT}的值不宜过低，否则会导致计算结果不准确
	\item \textrm{ENCUT}也不宜过高，否则会过度消耗计算资源，增加计算成本
	\item 一般建议\textrm{ENCUT}的值设置为赝势\textrm{ENMAX}的1.0到1.3倍
\end{itemize}

\textrm{ALGO}:~确定电子步能量最小化算法。
\begin{itemize}
	\item \textrm{Normal}:~标准的\textrm{Davidson}块迭代对角化算法，较为稳定可靠(属于默认值)
	\item \textrm{Fast}:~速度更快，但可能不够稳定
	\item \textrm{All}:~使用完全共轭梯度算法\textrm{(Conjugate Gradient, CG)}
\end{itemize}

\textrm{NELM}:~指定电子自洽迭代的步数上限。默认值为60，\textcolor{red}{若计算未能在该步数内收敛，可适当增大}

\textrm{EDIFF}:~电子自洽迭代的收敛标准。当两次迭代之间的总能量变化小于\textrm{EDIFF}，电子迭代收敛。默认值通常为\textrm{1E-4($10^{-4}$)}或$10^{-5}$，对于高精度计算，建议设为$10^{-6}$或更小

\textrm{ISMEAR}与\textrm{SIGMA}:~确定刻画导体\textrm{Fermi}电子占据数的展宽方案和展宽大小
\begin{itemize}
	\item \textrm{ISMEAR = -5}:~四面体\textrm{(Tetrahedron)}方案。%适用于DOS计算和绝缘体/半导体的精确总能计算，不适用于金属的结构弛豫）。
	\item \textrm{ISMEAR = 0}:~\textrm{Gaussian}展宽方案
	\item \textrm{ISMEAR = -1}:~\textrm{Fermi}展宽方案
	\item \textrm{ISMEAR = 1}:~\textrm{Methfessel-Paxton}方案，展宽逼近取到一阶%（适用于金属）。
	\item \textrm{SIGMA}~展宽宽度参数，通常设置在\textrm{0.01$\sim$0.2 eV}之间
\end{itemize}

\subsubsection{离子弛豫与结构优化参数}
\textrm{IBRION}:~离子运动或弛豫算法
\begin{itemize}
	\item \textrm{-1}:~静态计算
	\item \textrm{0}:~分子动力学模拟
	\item \textrm{1}:~%\textrm{RMM-DIIS}
		算法\textrm{quasi-Newton}方案
	\item 共轭梯度算法:~\textcolor{red}{通常较稳定，结构优化优先推荐}
\end{itemize}

\textrm{ISIF}:~结构优化过程中许可变化的自由度控制参数。\textcolor{red}{\textrm{ISIF}是结构优化中最关键的参数之一}
\begin{itemize}
	\item 0:~不计算应力张量
	\item 1:~仅计算应力，不允许晶胞体积和形状发生变化
	\item 2:~仅允许离子位置发生变化，不改变晶胞体积和形状
	\item 3:~允许离子位置、晶胞形状和体积同时发生变化~(完全弛豫)
	\item 4:~固定晶胞体积，但允许改变晶胞形状和离子位置发生变化
	\item 7:~仅允许晶胞的体积发生变化
\end{itemize}

\textrm{NSW}:~计算过程中控制离子步数上限的参数。\textcolor{red}{当离子步数达到\textrm{NSW}时，计算停止。对于结构优化来说，离子弛豫通常设置为100或200。当NSW=0时，表示完成单点能(静态计算)计算}

\textrm{EDIFFG}:~设定离子弛豫收敛标准的参数
\begin{itemize}
	\item 若\textrm{EDIFFG>0}:~表示该值为能量收敛的阈值
	\item 若\textrm{EDIFFG<0}:~表示该值为最大原子受力收敛的阈值
\end{itemize}
通常建议\textrm{EDIFFG}采用负值(如\textrm{EDIFFG = -0.02}表示所有原子受力小于\textrm{0.02 eV/\AA}时收敛)，一般来说，力收敛比能量收敛标准更为严格。若\textrm{EDIFFG}未显式设置，其默认值一般取为$10\times\mathrm{EDIFF}$。

\textrm{POTIM}:~分子动力学计算中的离子步的初始时间步长，默认值通常取为0.5

\subsubsection{计算精度控制与输出参数}
\textrm{PREC}:~全局计算精度，可选值为\textrm{Low, Medium, Normal, Accurate, High}
\begin{itemize}
	\item 参数\textrm{PREC}值的选择，将会影响到\textrm{ENCUT}，\textrm{FFT}网格等多个相关参数的默认值
	\item 一般计算中，\textrm{Normal}或\textrm{Accurate}是最常见的选取值
\end{itemize}

\textrm{LREAL}:~控制是否实空间计算投影函数的参数
\begin{itemize}
	\item \textrm{.FALSE.}:~仅在倒空间计算投影函数\textcolor{red}{(对较小体系效率更高)}
	\item \textrm{.TRUE.}:~允许在实空间计算投影函数
	\item \textrm{Auto}:~由计算程序自动选择确定
\end{itemize}

\textrm{LORBIT}:~控制投影态密度\textrm{projected DOS}的输出参数
\begin{itemize}
	\item 10:~结果输出到\textrm{PROCAR}文件
	\item 11:~结果输出到\textrm{DOSCAR}和\textrm{PROCAR}
\end{itemize}

\textrm{LWAVE}:~控制输出\textrm{WAVECAR}文件。\textcolor{red}{对于较大体系，建议设为\textrm{.FALSE.}以节省存储空间}

\textrm{LCHARG}:~控制输出\textrm{CHGCAR}电荷密度文件参数

\subsection{INCAR}文件的参数值
完整的\textrm{INCAR}中的各参数，包括没有设置采用默认值的，都会在\textrm{OUTCAR}中给出。以下是一个典型的结构优化的全部\textrm{INCAR}参数(\textcolor{red}{以体相铜为例}):
%\textcolor{red}{\textbf{直接插入文件}}:
{\fontsize{8.2pt}{3.2pt}\selectfont{
\verbatiminput{Figures/OUTCAR-INPUT.txt} %为保险:~选用文件名绝对路径
}}
%```
%SYSTEM = Cu bulk optimization
%PREC = Normal
%ENCUT = 520
%EDIFF = 1E-6
%EDIFFG = -0.02
%IBRION = 2
%ISIF = 3
%NSW = 100
%ISMEAR = 1
%SIGMA = 0.2
%LREAL = Auto
%LORBIT = 11
%LWAVE = .FALSE.
%LCHARG = .TRUE.
%```
\textrm{INCAR}中的常见错误
\begin{enumerate}
	\item 参数拼写错误:~如将\textrm{SYSTEM}误写为\textrm{RYSTEM}，将会导致计算结果出现异常
	\item 参数冲突:~同时开启或指定互不兼容的任务类型~(如分子动力学计算与结构优化的参数混用)
	\item 能量截断过低:~计算使用的基组过小，可能导致结果不可靠
	\item 收敛标准过于低:~导致未收敛的结果被误认为已经达到收敛
	\item \textrm{ISIF}参数选择不当:~例如，在表面计算中，使用\textrm{ISIF=3}，结果导致晶胞在真空层方向也因弛豫而被优化
\end{enumerate}

\section{\rm{KPOINRS}文件}% ## 四、KPOINTS文件：布里渊区采样的“方案”
\subsection{$\vec k$空间布点与积分}
\textrm{KPOINTS}是\textrm{VASP}中定义$\vec k$点(\textrm{Brillouin zone}空间的布点)和采样方案的文件。在周期体系的电子结构计算中，需要对第一\textrm{Brillouin zone}的布点进行积分，而$\vec k$点采样就是在倒空间中对该积分进行数值离散化的方案。$\vec k$点网格的布点密度直接决定了计算的精度和计算量——一般地说，网格越密，计算结果越精确，但计算成本也越高。因此，理性地选择$\vec k$点网格是平衡精度与效率的重要手段之一。
\subsection{\rm{KPOINTS}文件格式}
\textrm{KPOINTS}文件是\textrm{VASP}的输入文件中格式最固定的一个，最常见的格式是五行模板:
\begin{itemize}
 \item 第1行:~注释，可以是任意字符串(但不能为空)，一般是$\vec k$空间布点的相关描述
 \item 第2行:~整数，0表示$\vec k$空间网格点将是通过自动生成
 \item 第3行:~$\vec k$空间网格分布类型\\
	 \textrm{G}或\textrm{Gamma}:~表示以$\Gamma$点为中心的$\vec k$点网格生成方案\\
	 \textrm{M}或\textrm{Monkhorst}:~表示$\vec k$空间网格将采用\textrm{Monkhorst-Pack}生成方案。
%两者的区别在于：Monkhorst-Pack网格在三个方向平移了1/(2N)个单位，而Gamma中心网格是Monkhorst-Pack网格的特殊情况（平移量为0）。对于六角晶格等体系，强烈推荐使用Gamma中心网格。

 \item 第4行:~三个整数，分别表示在倒格矢$\mathbf{b}_1$、$\mathbf{b}_2$、$\mathbf{b}_3$三个方向上的$\vec k$点数目\\
	 例如~4 4 4~表示在每个方向取4个$\vec k$点，共64个k点。

 \item 第5行:~三个实数，表示网格原点的偏离量\textrm{(shift)}，通常设为~0 0 0。\textcolor{red}{\textrm{Monkhorst-Pack}网格布点的中心相当于$\Gamma$布点偏移了0.5 0.5 0.5}
\end{itemize}

\textcolor{red}{$\vec k$空间网格布点示例}
{\fontsize{9.2pt}{7.2pt}\selectfont{
\verbatiminput{Figures/KPOINTS.txt} %为保险:~选用文件名绝对路径
}}
%Automatic generation
%0
%Gamma
%4 4 4
%0 0 0
\subsection{\rm{KPOINTS}的网格点测试选取}
\begin{itemize}
	\item 收敛性测试\\
		\textcolor{red}{最科学的方法是对目标体系进行k点收敛性测试。逐渐增加k点密度，观察总能量等关键物理量是否收敛，从而确定所需的最小k点数目}
	\item 经验规则\\
%- 对于块体材料，每个晶格矢量长度（Å）乘以该方向上的k点数目，通常在30~40之间。
%- 例如，晶格常数为6 Å的体系，k点可取3×3×3（6×3=18偏少）或5×5×5（6×5=30合适）。
%- 对于绝缘体，可以用较少的k点；对于金属，由于费米面附近需要更精细的采样，需要更多的k点。

	\item 使用\textrm{VASPKIT}
		\textrm{VASPKIT}工具可以自动生成\textrm{KPOINTS}文件\\
%		用户指定倒空间中k点间距（通常为0.04 Å⁻¹用于常规计算，0.03或0.02 Å⁻¹用于高精度计算），程序自动生成合适的k点网格。

%	\item INCAR自动生成（替代KPOINTS）**
%可以在INCAR中设置`KSPACING`和`KGAMMA`来让VASP自动确定k点网格。这种方法在VASP官网中也被推荐。
%- `KSPACING`：决定k点在倒空间的密度。
%- `KGAMMA`：决定是否强制包含Gamma点。
\end{itemize}

\subsection{一些特殊情况的\rm{KPOINTS}设置}
真空和表面体系\\
对于包含真空层的表面或二维材料计算，真空层方向的k点应取1。%例如表面计算常用`4 4 1`。

能带计算\\
能带结构计算需要使用沿高对称点方向的路径，即\textrm{Line-mode}%，而非均匀网格。
\textcolor{red}{\textbf{直接插入文件}}:
{\fontsize{9.2pt}{7.2pt}\selectfont{
\verbatiminput{Figures/KPOINTS-Line.txt} %为保险:~选用文件名绝对路径
}}
分子/原子体系\\
对于孤立的分子或原子，或使用较大的晶胞，通常只需要$\Gamma$点~$(1\times1\times1)$。

%\textrm{\rm{KPOINTS}的}
%1. **第2行误写为字母**：很多人会将第二行的数字`0`误写为字母`o`。
%2. **真空层方向取了多个k点**：在包含真空层的方向上取k点不仅没有物理意义，还浪费计算资源。
%3. **k点过密或过疏**：过密浪费机时，过疏导致结果不准确。
%4. **未做收敛性测试**：任何正式计算前都应做k点收敛性测试。
%
%\subsection{\rm{KPOINTS}文件的内容}
%\subsection{\rm{VASP}中的$\vec k$空间积分}
\section{\rm{POTCAR}文件}
习惯上，\textrm{POTCAR}被成为\textrm{VASP}计算的``赝势文件''，但严格地说，\textrm{POTCAR}不仅仅是包含了计算体系中各元素的赝势数据，还包含了原子的波函数、电荷密度和投影函数等\textrm{PAW}方法所必须的各类基础数据，所以称其为``原子数据集文件''可能更准确。而且\textrm{VASP}的相关代码都是开源的，但唯有\textrm{POTCAR}的生成方案至今没有公开，\textrm{VASP}提供的较为完整的\textrm{POTCAR}数据，是\textrm{VASP}计算结果可靠的重要前提。
\subsection{\rm{POTCAR}文件的内容}
以下结合铜\textrm{(Cu)}原子的\textrm{POTCAR}为示例，说明\textrm{POTCAR}文件的主要组成。\textrm{POTCAR}文件主要包含以下若干部分：
\begin{itemize}
	\item \textrm{POTCAR}数据生成的控制参数
{\fontsize{8.2pt}{6.2pt}\selectfont{
\verbatiminput{Figures/POTCAR-1-parameter.txt} %为保险:~选用文件名绝对路径
}}
字符串``PSCTR-controll parameters''为控制参数结束的标记。这些参数主要是原子名称、价电子构型、超软赝势\textrm{(USPP)}\textrm{PAW}方法产生原子赝势所需的截断半经等各类参数，以及原子态电子的动能误差信息。熟悉赝势和\textrm{PAW}方法的读者不难猜出，这些参数应该是用于生成\textrm{POTCAR}原子数据的程序\textrm{PSCTR}(\textrm{VASP}未公开的代码)所需的各类控制参数。
\item 原子赝势局域部分在倒空间的表示
{\fontsize{8.2pt}{6.2pt}\selectfont{
	\lstinputlisting[linerange=1-5]{Figures/POTCAR-2-local_G.txt} %为保险:~选用文件名绝对路径
}}
$\cdots\cdots$
{\fontsize{8.2pt}{6.2pt}\selectfont{
	\lstinputlisting[linerange=199-204]{Figures/POTCAR-2-local_G.txt} %为保险:~选用文件名绝对路径
}}
以字符串``local part''为起始标记，开始\textrm{USPP}和\textrm{PAW}都会涉及到的原子赝势(局域部分)。紧随其后的是倒空间中取$\vec G$的截断值；随后为$V_{\mathrm{loc}}$的1000个数值点，即对应将$\vec G$值均分1000份后，$V_{\mathrm{loc}}(G)$的各数值点。此后的字符串``gradient corrections used for XC''提示交换-相关泛函，其随后的数值5表示该处原子势中的交换-相关部分采用\textrm{PBE}泛函形式。
 \item 原子赝电荷密度的局域部分在倒空间的表示
{\fontsize{8.2pt}{6.2pt}\selectfont{
	\lstinputlisting[linerange=1-5]{Figures/POTCAR-3-charge_G.txt} %为保险:~选用文件名绝对路径
}}
$\cdots\cdots$
{\fontsize{8.2pt}{6.2pt}\selectfont{
	\lstinputlisting[linerange=199-204]{Figures/POTCAR-3-charge_G.txt} %为保险:~选用文件名绝对路径
}}
$\cdots\cdots$
{\fontsize{8.2pt}{6.2pt}\selectfont{
	\lstinputlisting[linerange=399-402]{Figures/POTCAR-3-charge_G.txt} %为保险:~选用文件名绝对路径
}}
原子价电子密度的局域部分以字符串``core charge-density (partial)''为起始标记，紧随其后为$n_{\mathrm{loc}}^{\mathrm{core}}(G)$的1000个数值点；以字符串``atomic pseudo charge-density''为赝电荷密度的起始标记，随后是$\tilde n_{\mathrm{loc}}^{\mathrm{core}}(G)$的1000个数值点。
\item 投影子函数在倒空间和实空间的表示
{\fontsize{8.2pt}{6.2pt}\selectfont{
	\lstinputlisting[linerange=1-8]{Figures/POTCAR-4-projector.txt} %为保险:~选用文件名绝对路径
}}
$\cdots\cdots$
{\fontsize{8.2pt}{6.2pt}\selectfont{
	\lstinputlisting[linerange=24-30]{Figures/POTCAR-4-projector.txt} %为保险:~选用文件名绝对路径
}}
$\cdots\cdots$
{\fontsize{8.2pt}{6.2pt}\selectfont{
	\lstinputlisting[linerange=44-50]{Figures/POTCAR-4-projector.txt} %为保险:~选用文件名绝对路径
}}
熟悉赝势理论的读者都知道，投影子函数$\tilde p_l$与原子波函数相关，其构造依赖角动量$l$和赝波函数的能量。与势函数的$\vec G$空间布点数(1000)不同，投影子函数在$\vec k$空间和实空间的布点数都为100。因在倒空间的截断值列于第一行，并跟随字符``\textrm{T}''标注。其下字符串``Non local Part''标记非局域部分的投影子函数相关的贡献，紧随标记行的三个参数分别为投影子函数的轨道角动量$l$、轨道数目和实空间径向截断半径。再下一行则是原子超软赝势或\textrm{PAW}方法所涉及的非局域部分贡献的$D^{\mathrm{ion}}$项(共四个)，这几个数据的获得，恰恰是\textrm{USPP}或\textrm{PAW}方法的有价值部分。随后即为分别用字符串``Reciprocal Space Part''和``Real Space Part''标记的轨道角动量$l$的投影子函数在倒空间和正空间的数值。具体到\textrm{Cu}元素，依次为$l=2, 0, 1$的各两组投影子函数。
\item 基于\textrm{PAW}方法的参数
{\fontsize{8.2pt}{6.2pt}\selectfont{
	\lstinputlisting[linerange=1-8]{Figures/POTCAR-5-PAW-grid.txt} %为保险:~选用文件名绝对路径
}}
$\cdots\cdots$
{\fontsize{8.2pt}{6.2pt}\selectfont{
	\lstinputlisting[linerange=10-15]{Figures/POTCAR-5-PAW-grid.txt} %为保险:~选用文件名绝对路径
}}
$\cdots\cdots$
{\fontsize{8.2pt}{6.2pt}\selectfont{
	\lstinputlisting[linerange=20-24]{Figures/POTCAR-5-PAW-grid.txt} %为保险:~选用文件名绝对路径
}}
$\cdots\cdots$
{\fontsize{7.2pt}{5.2pt}\selectfont{
	\lstinputlisting[linerange=79-86]{Figures/POTCAR-5-PAW-grid.txt} %为保险:~选用文件名绝对路径
}}
\textrm{PAW}方法的参数以字符串``\textrm{PAW radial sets}''起始作标记。紧随标记字符串后是指定的径向(按对数布点)的布点数和``截断半径''(真实采用的``截断半径''会比该数值略大)；紧随后一行字符串则是用作指定后续实数型数据的格式，可能是突出\textrm{PAW}方法所用的数据格式特征。再次行的字符串``\textrm{augmentation charges (non sperical)}''起始标记的是赝势中的补偿电荷。字符串``\textrm{uccopancies in atom}''起始标记的是原子价电子轨道的占据情况，有必要指出的是，这里占据电子是平均分配到每个轨道的，因此此处的电子数会存在分数占据。字符串``\textrm{grid}''起始标记的是径向对数布点的各点数值。在\textrm{Cu}的实例中，指定``截断半径''对应第313点(倒数第七点)。
\item 原子全电子势函数
{\fontsize{7.2pt}{5.2pt}\selectfont{
	\lstinputlisting[linerange=1-5]{Figures/POTCAR-5-aepotential.txt} %为保险:~选用文件名绝对路径
}}
$\cdots\cdots$
{\fontsize{7.2pt}{5.2pt}\selectfont{
	\lstinputlisting[linerange=63-65]{Figures/POTCAR-5-aepotential.txt} %为保险:~选用文件名绝对路径
}}
自此以下数据与标注的关系都较为清晰明确，主要是\textrm{PAW}方法的原子数据。原子的全电子势函数以字符串``aepotential''起始标记。
\item 原子芯电荷密度
{\fontsize{7.2pt}{5.2pt}\selectfont{
	\lstinputlisting[linerange=1-5]{Figures/POTCAR-6-core.txt} %为保险:~选用文件名绝对路径
}}
$\cdots\cdots$
{\fontsize{7.2pt}{5.2pt}\selectfont{
	\lstinputlisting[linerange=63-65]{Figures/POTCAR-6-core.txt} %为保险:~选用文件名绝对路径
}}
类似地，字符串``\textrm{core charge-density}''起始标记的芯电荷密度。  
\item 原子动能密度
{\fontsize{7.2pt}{5.2pt}\selectfont{
	\lstinputlisting[linerange=1-5]{Figures/POTCAR-7-kinetic.txt} %为保险:~选用文件名绝对路径
}}
$\cdots\cdots$
{\fontsize{7.2pt}{5.2pt}\selectfont{
	\lstinputlisting[linerange=63-65]{Figures/POTCAR-7-kinetic.txt} %为保险:~选用文件名绝对路径
}}
字符串``\textrm{kinetic energy-density}''起始标记的芯电荷密度。  
\item 原子的赝势
{\fontsize{7.2pt}{5.2pt}\selectfont{
	\lstinputlisting[linerange=1-5]{Figures/POTCAR-8-pspotential.txt} %为保险:~选用文件名绝对路径
}}
$\cdots\cdots$
{\fontsize{7.2pt}{5.2pt}\selectfont{
	\lstinputlisting[linerange=63-65]{Figures/POTCAR-8-pspotential.txt} %为保险:~选用文件名绝对路径
}}
字符串``pspotential''起始标记原子的赝势。  
\item 原子的赝芯电荷密度
{\fontsize{7.2pt}{5.2pt}\selectfont{
	\lstinputlisting[linerange=1-5]{Figures/POTCAR-8-core-ps.txt} %为保险:~选用文件名绝对路径
}}
$\cdots\cdots$
{\fontsize{7.2pt}{5.2pt}\selectfont{
	\lstinputlisting[linerange=63-65]{Figures/POTCAR-8-core-ps.txt} %为保险:~选用文件名绝对路径
}}
字符串``\textrm{core charge-density (pseudized)}''起始标记原子的赝势。
\item 原子的价电子的赝波函数和全电子波函数
{\fontsize{7.2pt}{5.2pt}\selectfont{
	\lstinputlisting[linerange=1-5]{Figures/POTCAR-9-wave-6.txt} %为保险:~选用文件名绝对路径
}}
$\cdots\cdots$
{\fontsize{7.2pt}{5.2pt}\selectfont{
	\lstinputlisting[linerange=63-68]{Figures/POTCAR-9-wave-6.txt} %为保险:~选用文件名绝对路径
}}
$\cdots\cdots$
{\fontsize{7.2pt}{5.2pt}\selectfont{
	\lstinputlisting[linerange=128-131]{Figures/POTCAR-9-wave-6.txt} %为保险:~选用文件名绝对路径
}}
不同角动量$l$的价电子波函数与投影函数的排序严格对应。很显然，字符串``\textrm{pseudo wavefunction}''起始标记原子的赝波函数；字符串``\textrm{ae wavefunction}''起始标记原子的全电子波函数。具体到\textrm{Cu}元素，原子波函数依次为$l=2, 0, 1$的各两组轨道函数。
\end{itemize}
\textrm{POTCAR}的最后，原子数据集以字符串``End of Dataset''标注结束。
%\subsection{\rm{POTCAR}文件与\rm{PAW}方法}
%赝势方法的核心思想是将原子的内层电子（芯电子）与原子核合并为一个有效的“离子实”，用一个平滑的势函数来代替真实的原子势，从而大幅减少计算所需的平面波数量。POTCAR文件告诉VASP计算体系中包含哪些元素以及使用哪种赝势来描述这些元素。
%## 二、POSCAR文件：原子结构的“骨架”
%
%### 2.1 POSCAR文件概述
%
%POSCAR是VASP中描述原子结构的输入文件。它包含了计算体系的所有结构信息——晶胞的晶格参数、原子种类与数量、各原子的空间位置，以及可选的原子初始速度（用于分子动力学模拟）。POSCAR是一个纯文本文件，每一行都有特定的格式含义。
%
%VASP在计算结束后会自动生成一个CONTCAR文件，其格式与POSCAR几乎完全相同，记录了计算后体系的最终几何结构。CONTCAR可以直接重命名为POSCAR用于后续计算（如更精确的静态计算或性质计算）。
%
%### 2.2 POSCAR文件格式详解
%
%一个标准的POSCAR文件通常包含以下行（按顺序）：
%
%**第1行：标题行（注释行）**
%任意字符串，用于描述体系（如“Si bulk”或“FePt alloy”），不影响计算，仅用于识别。
%
%**第2行：全局缩放因子**
%一个标量值，用于统一缩放后续所有晶格向量。通常设置为1.0。若设为0.5，则所有晶格向量缩小一半。
%
%**第3-5行：晶格向量（3行）**
%定义晶胞的三个基矢a、b、c，每行包含三个浮点数，单位为埃（Å）。这三行向量定义了晶胞的形状和大小。
%
%**第6行：原子种类**
%列出体系中包含的元素符号（如“Si O”或“Fe Pt”）。该行可省略，若省略则需要在第7行直接给出原子数量。
%
%**第7行：各元素原子数**
%按第6行元素顺序，列出每种元素的原子个数。例如“2 4”表示第一种元素有2个原子，第二种有4个原子。
%
%**第8行：坐标类型**
%- `Direct`或`D`：分数坐标（坐标值在0~1之间，相对于晶格向量）。
%- `Cartesian`或`C`：笛卡尔坐标（绝对坐标，单位为Å）。
%- Direct坐标是最常用的选择，因为它与晶格参数直接关联，便于晶胞形变时的坐标转换。
%
%**第9行及之后：原子坐标**
%每个原子占一行，按第6-7行指定的顺序排列。若使用Direct坐标，每行三个数均在0~1之间；若使用Cartesian坐标，则为以Å为单位的绝对坐标。
%
%### 2.3 选择性动力学（Selective Dynamics）
%
%在POSCAR中，可以在坐标类型行之前插入一行“Selective dynamics”（或仅以字母S开头）来启用选择性动力学模式。启用后，每个原子坐标后面可以附加三个逻辑标志（T或F），分别控制该原子在x、y、z三个方向上的移动是否被允许。
%
%例如，固定某个原子的所有自由度：
%```
%Selective dynamics
%Direct
%0.00 0.00 0.00 F F F
%```
%这表示该原子在结构优化过程中位置完全固定。
%
%选择性动力学在以下场景中特别有用：
%- **表面计算**：固定底层原子，弛豫表层原子。
%- **吸附计算**：固定基底原子，只弛豫吸附分子。
%- **特定方向优化**：例如仅允许原子在z方向移动。
%
%### 2.4 POSCAR文件示例
%
%**示例1：体相硅（Si），Direct坐标**
%```
%Si bulk
%1.0
%0.0 2.715 2.715
%2.715 0.0 2.715
%2.715 2.715 0.0
%Si
%2
%Direct
%0.00 0.00 0.00
%0.25 0.25 0.25
%```
%
%**示例2：表面体系（含选择性动力学）**
%```
%Ni(100) surface
%1.0
%3.53 0.00 0.00
%0.00 3.53 0.00
%0.00 0.00 20.00
%Ni
%12
%Selective dynamics
%Direct
%0.00 0.00 0.00 F F F
%0.50 0.50 0.00 F F F
%0.00 0.50 0.10 T T T
%...
%```
%
%### 2.5 POSCAR常见注意事项
%
%1. **原子顺序一致性**：POSCAR中元素的排列顺序必须与POTCAR中赝势文件的排列顺序完全一致。
%2. **坐标类型选择**：Direct坐标最常用，因为它与晶格参数直接关联。
%3. **真空层设置**：表面计算时，沿非周期方向（如z轴）需添加足够的真空层（通常≥15 Å）以避免层间镜像相互作用。
%4. **缩放因子**：通常保持为1.0，避免引入额外误差。
%
%## 三、POTCAR文件：赝势的“能量密码本”

%### 3.1 POTCAR文件概述
对用户来说，\textrm{VASP}提供了较为完善的赝势库，用户所需要做的，是选择适合自己计算体系的原子数据集类型，将相应的原子数据文件组合成最终计算所需要的\textrm{POTCAR}文件。

%### 3.2 赝势的分类
%
%VASP的赝势可以从两个维度进行分类：
%
%**（一）按产生方法分类**
%- **USPP（Ultrasoft Pseudopotential，超软赝势）** ：较早发展的赝势方法，通过减小截断能来提高计算效率。
%- **PAW（Projector Augmented Wave，缀加平面波）** ：更先进的方法，在保留赝势效率的同时，能够重建芯区域的真实价电子波函数。
%
%VASP官方强烈推荐使用PAW赝势。原因有二：一是PAW的径向截断半径（芯半径）比USPP更小，因此更精确；二是PAW方法能重构芯区域带有节点的真实价电子波函数。USPP可以看作是早期计算能力受限时的一种妥协方案，如今已很少使用。

%- **LDA（Local Density Approximation，局域密度近似）** 。
%- **GGA（Generalized Gradient Approximation，广义梯度近似）** ，其下
%目前最常用的是**PAW + PBE**组合。

\subsection{\rm{VASP}原子数据集的结构}
\textcolor{red}{在发布的\textrm{VASP}软件中，原子数据集文件通常存放在`$\sim$/vasp/potentials`目录下，}习惯上，这些原子数据集，按交换-相关泛函的不同，划分为\textrm{LDA}型和\textrm{GGA}型，\textrm{GGA}型又可再细分为\textrm{PW91}和\textrm{PBE}型等。因此原子数据集的目录，主要包含以下几个子目录:

\begin{table}[h!]
    \centering
    \begin{tabular*}{\textwidth}{|@{\extracolsep{\fill}}c|@{\extracolsep{\fill}}c|@{\extracolsep{\fill}}c|}
        \toprule
目录 & 赝波函数和赝势生成方案 & 交换-相关泛函\\
\midrule
\textrm{pot} & \textrm{PP}~(\textrm{USPP}为主) & \textrm{LDA}\\
\textrm{pot\_GGA} & \textrm{PP}~(\textrm{USPP}为主) & \textrm{GGA}\\
\textrm{potpaw} & \textrm{PAW} & \textrm{LDA}\\
\textrm{potpaw\_GGA} & \textrm{PAW} & \textrm{GGA~(PW91)}\\
\textrm{potpaw\_PBE} & \textrm{PAW} & \textrm{GGA~(PBE)}\\
        \bottomrule
    \end{tabular*}
    \caption{\textrm{POTCAR}的类型组织。}
\label{Tab:POTCAR-Tag-0}
\end{table}

在每个元素的目录下，常常还会将原子数据集进一步划分。以元素\textrm{Ga}为例，共有
\begin{itemize}
	\item \textrm{Ga}最通用的原子数据集
	\item \textrm{Ga\_s}:~\textrm{s}表示\textrm{softer}，产生数据时的截断能较低。计算较快但精度略低
	\item \textrm{Ga\_h}:~\textrm{s}表示\textrm{harder}，产生数据时的截断能较高。计算精度较高高
	\item \textrm{Ga\_sv}:~将接近价层的芯层\textit{s}电子(半芯层)，作为价电子处理
	\item \textrm{Ga\_pv}:~将接近价层的芯层\textit{p}电子(半芯态)，作为价电子处理
\end{itemize}
%如果关注的是d电子的精确描述，可能需要选择包含半芯态的赝势（如_sv）。

\subsection{计算用\rm{POTCAR}文件的生成}
计算用的\textrm{POTCAR}文件，注意将所需各元素的原子数据集文件，按\textrm{POSCAR}中原子的顺序，依次合并。以\ch{GaAs}为例，如果\textrm{POSCAR}中是\textrm{Ga As}，则\textrm{POTCAR}中的原子数据集也必须是\textrm{Ga}在前、\textrm{As}在后。

%1. 进入赝势目录（如`potpaw_PBE`），找到Ga和As的赝势文件。
%2. 解压并复制到工作目录：
%   ```bash
 %  zcat POTCAR.Z > Ga
 %  zcat POTCAR.Z > As
 %  ```

%3. 按顺序合并：
%   ```bash
%   cat Ga As > POTCAR
%   ```
同一计算中，所有元素的赝势建议来自同一类型(如同为\textrm{PAW\_PBE})，不宜混用不同类型的原子数据。

%**关键提醒**：POTCAR中元素的排列顺序必须与POSCAR中元素的排列顺序完全一致。例如，
%
%### 3.5 POTCAR文件中的关键信息
%
%POTCAR文件包含许多对计算有重要影响的信息，可以通过以下命令查看：
%
%```bash
%grep -in "enmax" POTCAR
%```
%
%这会显示POTCAR中推荐的截断能（ENMAX），该值通常被用作INCAR中ENCUT的参考基准。实际计算中，ENCUT通常设置为ENMAX的1.0到1.3倍。
%
%### 3.6 赝势选择的实践建议
%
%1. **首选PAW_PBE**：对于绝大多数体系，PAW方法结合PBE泛函是最稳妥的选择。
%2. **注意元素特异性**：对于过渡金属等需要精确描述d电子或f电子的体系，可能需要选择包含半芯态的赝势（如_sv或_pv）。
%3. **保持一致性**：同一计算中所有元素的赝势应来自同一族（如同为PAW_PBE），不可混用不同版本。
%4. **做收敛性测试**：对于不熟悉的体系，建议测试不同赝势对计算结果的影响。

%## 五、四个输入文件的协同与关联
%
%四个输入文件虽然在功能上各司其职，但彼此之间存在紧密的关联：
%
%**1. POSCAR与POTCAR的元素顺序必须一致**
%这是VASP计算中最基本也是最重要的规则之一。若POSCAR中元素顺序为`Si O`，则POTCAR必须是`Si`的赝势在前、`O`的赝势在后。
%
%**2. INCAR中的ENCUT与POTCAR中的ENMAX密切相关**
%ENCUT应至少等于POTCAR中各元素ENMAX的最大值，通常建议设为最大ENMAX的1.0~1.3倍。
%
%**3. POSCAR的晶格参数与KPOINTS的网格密度相互影响**
%晶胞越大，在倒空间中所需的k点越少（因为倒格矢更短）；反之，晶胞越小，所需的k点越多。这就是所谓的“a×k ≈ 常数”经验规则。
%
%**4. INCAR中的ISIF与POSCAR中的选择性动力学共同控制结构优化**
%ISIF决定哪些宏观自由度（晶胞体积、形状）可以变化，而选择性动力学控制哪些原子的哪些方向可以移动。两者配合使用可以实现精细的结构优化控制。
%
%**5. INCAR中的ISMEAR/SIGMA与KPOINTS的密度共同决定金属计算的精度**
%对于金属体系，需要足够密的k点网格和适当的展宽方法（如ISMEAR=1或-1）才能正确描述费米面附近的电子占据。
%
%## 六、输入文件的准备流程与最佳实践
%
%### 6.1 标准准备流程
%
%1. **确定计算体系**：明确要计算的物质种类、晶体结构、计算目的（结构优化、静态计算、DOS、能带等）。
%2. **准备POSCAR**：从实验数据、数据库或建模软件获取晶体结构，按格式编写POSCAR。
%3. **选择并生成POTCAR**：选择合适的赝势类型（推荐PAW_PBE），按POSCAR中的元素顺序合并生成POTCAR。
%4. **确定KPOINTS**：根据体系大小和对称性，选择合适的k点网格，必要时做收敛性测试。
%5. **编写INCAR**：根据计算目的设置相应参数，从模板出发，逐步完善。
%6. **提交计算**：将四个文件放入同一目录，提交作业。
%
%### 6.2 最佳实践建议
%
%1. **从简单开始**：先用较宽松的参数（如较少的k点、较低的ENCUT、较宽松的EDIFF）做快速测试，确认结构无误后再提高精度。
%2. **做收敛性测试**：对ENCUT、k点密度、EDIFF等关键参数进行系统的收敛性测试，确保结果的可靠性。
%3. **善用CONTCAR**：结构优化结束后，将CONTCAR重命名为POSCAR进行后续的静态计算或性质计算。
%4. **检查OUTCAR**：计算完成后，仔细检查OUTCAR文件，确认计算是否正常收敛，是否存在警告或错误信息。
%5. **保持文件整洁**：不同计算任务使用不同的工作目录，避免文件混乱。
%## 七、总结
%
%VASP的四个输入文件——INCAR、POSCAR、POTCAR和KPOINTS——共同构成了第一性原理计算的基础框架。INCAR作为“控制中心”，通过丰富的参数标签决定了计算“做什么”和“如何做”；POSCAR作为“结构蓝图”，精确描述了原子体系的几何构型；POTCAR作为“能量密码本”，提供了各元素的赝势信息；KPOINTS作为“采样方案”，定义了布里渊区的取样方式。四个文件缺一不可，协同工作，共同决定了VASP计算的精度、效率和可靠性。
%
%理解每个文件的格式规范、关键参数及其物理含义，是正确使用VASP进行第一性原理计算的基础。初学者应从标准模板出发，逐步深入理解各参数的意义，并通过不断的实践积累经验。正如VASP官方手册所建议的：当不确定某个参数时，使用默认值通常是一个明智的选择。但更重要的原则是——无论做什么计算，正式生产之前务必做充分的收敛性测试。只有这样，才能确保计算结果是具有发表价值的科学数据，而非毫无意义的数字噪音。
